Since the COVID-19 pandemic, the share of U.S. universities adopting test-optional admission policies has more than doubled, reaching 80\%. In place of standardized test scores such as the SAT or ACT, students now rely on high school GPA, personal statements, and extracurricular achievements to demonstrate academic potential. While these changes may enhance access for students from underrepresented backgrounds (\citet{Bennett2021}), standardized exams remain a stronger predictor of university success than GPA (\citet{Chetty2020}). Mitigation of barriers for applying to higher education 

This paper provides a comprehensive analysis of the beneficiaries of non-standardized testing, focusing on student demographics, prior academic achievement, school characteristics, and subject choices. During the COVID-19 pandemic, the British government canceled standardized in-person A-level exams and allowed schools to submit teacher-predicted grades in their place. In the absence of a unified grading procedure, teachers assigned final grades based on their personal knowledge of each student. This policy led to a 12-percentage point increase in the share of top grades awarded. Because university placements in the UK are largely determined by A-level results and universities are required to admit all students who meet their grade offers, grade inflation contributed to widespread oversubscription at many institutions.

Leveraging this policy shift, I examine how grade assignment patterns differed between standardized and non-standardized regimes. I decompose the sources of the changes in the overall grade distribution into a reassignment of grades within a subject course class within a school ('\emph{within-school effect}') and the extent of the unilateral shifts in students' grade levels within a school (\emph{between-school effect}). I examine the within-school effect by showing how the extent of the change in the assigned grade correlates with the student's academic ability, measured by a standardized test that students take at the age of 16, and also by the student's economic and social affluence. I also examine the association between larger increases in grade levels and demographic attributes of students such as race and gender. I empirically test whether the nonstandardized grades lead to changes in the inequality on educational attainments across students within their schools. 

I then test whether schools systematically differed in their tendency to assign higher grades based on their attributes. I examine the association of larger leniency towards assigning higher grades with the school type (private school or public schools), and also with their institutional quality. I also test whether the grading patterns differed between students taking different subjects, which I particularly categorize by the extent of quantitative reasoning in their content. Quantitative contents may prevent teachers to over-report their students grades as the past scores of the students from past quiz had already ranked students which facilitated the teachers to predict grades. I estimate the grading patterns while taking into account the student's academic and socio-economic background, as their attending institutions tendency to assign higher grades as well as their sorting behavior.

Lastly, I provide predictions on how the two grading schemes lead to a different composition of students across universities. In the UK, universities control the number of students they admit by setting a uniform grade requirement that applicants must pass to secure a spot. Using a structural model with endogenous university application choices, with fixed number of seats at universities, I examine how the correlation between university rank and mean demographics of students differs across standardized and non-standardized exams. I document how the share of students from socioeconomically disadvantaged backgrounds at prestigious university institutions differ across the grading regimes. The exercise assesses the extent of the grade changes triggered by the regime change into improvements into placements into better universities, as well as the demographics of students that improves their placement outcomes and those that descend in their placement rank.

\vspace{1em}

I estimate the changes in grades assignments using a generalized mixed-effects model that distinguishes between student attributes that are common across all schools (e.g., prior academic performance, student demographics) and those specific to schools (e.g. school fixed effects). I measure the changes in grades by the tendency for a given student to receive a top grade. The model estimates the increments in the probability that a student will receive a top grade by comparing students in the nonpandemic years to students in previous cohorts with similar academic standings. To isolate the concentration of grade shifts for each attribute of students, the treatment variable is interacted with a rich set of student demographic measurements, as well as the school indicator. I also include those controls into the regression so that students in the pandemic years are matched with students who took the standardized exam in the nonpandemic years with academically and demographically similar attributes. 

I identify the treatment effects through an interrupted time series (ITS) design that compares the probability of receiving an A or A* across cohorts before and after the grading policy change. The concerns for identification are 1. The treatment effect may include underlying trends in the overall grade distribution, and 2. The treatment effect may not distinguish between changes in the grading method and disruption in student learning due to the closure of schools caused by the pandemic. For the first point, I exploit the institutional grading rule in the UK, in which the overall A-level grades followed a fixed distribution during the years leading to the pandemic. Examiners of the A level set grade boundaries to regulate grade inflation. I tested this property of the A-level distribution through multiple stationarity tests, showing that the average probability of receiving top grades remains stable over years when standardized grading was applied. For the second point, I provide evidence that the learning loss for the 2020 year cohort was small, as school closures were announced a month before the national examination date, and teachers were already informed about their students' academic competence through their teaching experience during the 2 years of the student learning period. [I provide national statistics of mock exams that were taken during the short span between the cancellation of the exams and the teachers' grading date, and show the non-existing evidence of learning disruption for the 2020 year cohort.]

I summarize the within-/between effect using an Oaxaca-Blinder-type decomposition that separates the two effects. I decompose the changes in the mean levels of students' chances of receiving a top grade into within-school effect, which keeps contribution of the school fixed effect as their baseline levels but allowing the student characteristic related factors to change their contribution (e.g. race, gender, previous academic achievements, etc.). I calculate the effect between schools similarly by shutting down the contribution of student characteristics. I also decompose the variance of the individual treatment effect, and decompose the variance into the extent of heterogeneity in both within/between school effects. 

To connect the grade changes to shifts into different university tiers, I develop a structural model of student application and university admission. The model estimates student preferences on universities while controlling for the universities' selectivity levels and the students' aggregate score of admission grades. The models maps grade distributions into the applicant pool for each university through the students' application portfolio. Through a styled assignment rule, in which spots in universities are taken by students from highest achievers, I estimate the full distribution of students across universities in the UK. I provide how the correlation between university selectivity levels and mean levels of the incoming student demographics, such as race and gender, differ between the standardized and non-standardized exam regimes.


\vspace{1em}

I find that within a subject course, the assignment of top grades was concentrated among students with academically and socioeconomic advantaged backgrounds. Students in the least economically deprived regions were 3 percentage points more likely to receive an A or A* compared to students in the most deprived regions. In addition, top grades were more likely to be assigned to students with higher scores in the previous academic tests in both mathematics and English. Students with lower academic or socio-economic backgrounds did not increase their chances of receiving a top grade. My results indicate that teachers predicted grades did not change the ordering of students final grade across students with varying baseline chances of receiving a top grade, but actually amplifying the inequality among students in the same class rooms to receive a top grade. 

Furthermore, I find that gender affected the chances of receiving a high grade after controlling for the students' previous academic performance. Female students were more likely to receive top grades than male students under the non-standardized grading regime. Among students with a baseline probability 50\% of receiving an A or A*, male students were 6 percentage points less likely to receive the highest grades than their female counterparts. This estimate remains robust to the inclusion of subject-level controls, suggesting that female students received higher grades than male students with comparable prior academic performance within the same classroom. In contrast, I find only weak differences in outcomes between white and non-white students: white students were awarded top grades at a rate 1 percentage point higher than non-white students, but this difference becomes statistically insignificant once school-level controls are included.

Students in socioeconomically more advantaged areas were significantly more likely to receive top grades. Furthermore, students from areas with a higher than average progression rate to higher education (64\%) had a 4-point higher probability of receiving top grades than students from regions with lower progression rates, averaging 18.6\%. These patterns suggest that residential context and local educational culture influenced grade assignment under the non-standardized regime.

I find that schools reduced their contributions to their students to receive a high grade. On average, schools reduced the probability for their average students to receive a higher grade by 2 percentage points. However, I find that the extent to which they reduced their contribution varied by school quality and their types. School with better quality, measured by their past cohorts exam performances at the A level, were associated with a larger degree of reducing their student's chances of receiving top grades, whereas schools with lower tendency to produce a top grade scoring student were associated with a lower degree in reducing their average students' chances of receiving a top grade. My results indicate that the nonstandardized exam reduced the overall contribution of the schools on students regressively, therefore flattening the educational inputs across schools in the standardized regime.

I also find varying grade inflation between school types, with private institutions exhibiting the largest upward shifts. Students in independent schools were 7 percentage points more likely to receive top grades than those attending three-year secondary schools (i.e., sixth-form colleges), which experienced the smallest increases. Independent schools were followed by six-year secondary schools (that is, state schools and academies), and then by selective public schools (i.e., grammar schools). In particular, the rise in top grades in private institutions was concentrated in 2021. In 2020, when university applications were submitted before teachers assigned grades, private school students experienced smaller increases in top-grade probabilities compared to students in universities and state schools. However, in 2021, when the cancellation of standardized exams occurred prior to the application deadline, students in private schools were over 5 percentage points more likely to receive top grades than their peers in public institutions. Finally, I find that schools with lower value-added measures exhibited greater increases in the probability that their students received top grades, suggesting that weaker academic input may have been offset by more generous grade assignment.

Lastly, I find that nonfacilitating subjects were associated with a higher probability of receiving top grades than facilitating subjects. On average, students who took nonfacilitating subjects were 8 percentage points more likely to receive an A or A* than those who took facilitating subjects. This difference persists even after controlling for the quantitative content of the subject, although the gap narrows to 6 percentage points. I also document evidence of intertemporal spillovers in grade inflation within subject groups. A one percentage point increase in the probability of receiving top grades in 2020 led to a 0.66 percentage point increase for facilitating subject takers and a 1.1 percentage point increase for non-facilitating subject takers in 2021. Similarly, inflation within subjects from 2020 to 2021 was 0.45 percentage points for facilitating subjects and 0.79 percentage points for non-facilitating subjects. These patterns suggest that grade inflation in one cohort may influence grading decisions in subsequent years, particularly in subjects where discretion is greater.

My structural model predicts that the income inequality across students across university worsens after the shift from standardized exam to non-standardized exam. The linear coefficient between the average parental income scores and the selectivity recorded of a university increases more than 5 times after replacement. The coefficient for the share of female students and the selectivity of the university decreases from -0.07 to -0.25, and the coefficient for the average math score decreases from -1.0 to -1.6. My estimate predicts how the nonstandardized grades lead to wealthier students to move up the university rank and replacing students with less economic affluence. 

My results describe how the removal of standardized exams led to a unification of achievement levels inputs across schools. Low quality schools took advantage of the grading discretion by employing a looser grading policy than high quality schools. In the resulting grade distribution, the gap in probabilities for the median student to receive a top grade in A-levels was reduced by [] percentage points. This aligns with \cite{figlio2004}'s findings of stronger concentration of grade manipulation at lower quality schools. 

I find that academically advantaged students increase their chances of receiving a top grade more than students with lower academic standings. This result aligns with the findings in \cite{TYNER2020102037}, that grade inflation tends to be larger for students with higher academic standings, regardless of their school of attendance. Although this suggests teachers or schools do not ignore the previous academic standings of their students, and potentially preserving the order of students in their classrooms by academic standings. However, my results suggest demographic backgrounds of students do contribute to differential grade inflation, as female students received higher grades than male students with similar academic standings. This aligns with \cite{lavy2018}'s findings, and suggests that discretionary grading benefits girls more as girls perform better in the daily classroom than boys, while under performing in high-pressure settings, such as the exam date. I do not find strong evidence of grade inflation by socioeconomic standings measured by the income deprivation index or regional proficiency to enter higher education. The findings are at odds with \cite{hanna2012} and \cite{burgess2013} which finds a positive relationship between socio-economic backgrounds and the extent of grade inflation.

High-quality schools tightened their grading criteria on their average students, as their high performing students received top grades due to their baseline chances of receiving a top grade being close to the boundary of receiving or not. The jump in the number of top-grade achievers compresses the grade increase for the students in the lower half of the performance distribution. On the other hand, students at lower quality schools increased their chances of receiving a top grade, as teachers may have the extra incentive to help an academically well performing but attending a bad school to place themselves at higher ranked academic institutions. 

However, the results do not necessarily translate to an increase in placements for students attending school with higher levels of grade shift. Although the regressive patterns of overall grade increases and school quality provides students attending low quality schools to obtain top grade, such students often do not apply to top ranked universities albeit receiving a high grade in A-levels. On the otherhand, students attending middle quality schools with good academic and socio-economic background leverage the upward shift in grades to secure placements at better universities. The upward shift of students from middle quality schools pulls students attending higher quality schools but with lower academic competence to lower their university placement. Despite positive relation between previous academic performance and the grade shift, the average academic score, measured by GCSE math, reduces at higher ranked universities. This indicates how non-academically related attributes, such as gender and socio-economic background dominates the upwards shifts due to academic performance. 

This paper contributes to the discussion on test-optional policies for university admissions. The results indicate how a removal of standardized exams changes the allocation of grades across the student population. 